% !PRINTED 312
\weberpage{312}
%% [page furniture, not body text -- commented out 2026-08-10] Page furniture (running
%% head, folio, signature) is excluded from this edition by standing policy.
%% \noindent Zehnter Abschnitt. \hfill \S.\,81.

If in the sum $\Sigma \alpha\xi$ we make all seven possible assumptions for the system $(\xi)$, then one obtains the value 0 three times and consequently the value 1 four times.

The congruence
\weq{7}{\Sigma \alpha\xi \equiv 1 \ (\mathrm{mod}\ 2)}
therefore has four and only four different solutions for $\xi$. The totality of these four solutions we shall call a \gesperrt{quadruple}. Every triple determines a quadruple uniquely and conversely, and so there are also seven different quadruples, which by (7) are completely determined by the number system $(\alpha)$. We now prove the theorem:

\gesperrt{6. If $(\xi)$ runs through the quadruple belonging to $(\alpha)$, and $(\beta)$ is a system different from $(\alpha)$, then among the four sums $\Sigma \beta\xi$ two are even and two are odd.}

For if $(\xi)$ runs through the totality of all seven systems, then among the sums $\Sigma \beta\xi$ there are found three even and four odd; of these, the triple belonging to $(\alpha)$ furnishes two odd and one even; there thus remain for the quadruple two even and two odd.

We now change the notation for the quadruples and take for the system $(\alpha)$ another symbol $(h_1, h_2, h_3) = (h)$. The quadruple belonging to $(h)$ we denote by $(\alpha), (\beta), (\gamma), (\delta)$, so that the four congruences:
\weq{8}{\Sigma h\alpha \equiv 1, \quad \Sigma h\beta \equiv 1, \quad \Sigma h\gamma \equiv 1, \quad \Sigma h\delta \equiv 1 \quad (\mathrm{mod}\ 2)}
hold.

In the denominator of the product of the four functions
\[
v_\alpha, \quad v_\beta, \quad v_\gamma, \quad v_\delta
\]
there occur by (6) in all 12 factors $\Psi_\xi$, and indeed each factor $\Psi_\xi$ occurs among them as often as the number of even numbers among the sums
\[
\Sigma \xi\alpha, \quad \Sigma \xi\beta, \quad \Sigma \xi\gamma, \quad \Sigma \xi\delta
\]
amounts to, i.\,e.\ $\Psi_h$ does not occur at all, while every other $\Psi_\xi$ occurs twice (by 6.).

If we now also make use of the relation (1), it follows that:
\weq{9}{\Psi_h^2 = A^2 v_\alpha v_\beta v_\gamma v_\delta.}
